\section{A finite-scale Gaussian result with no prior-strength limit}
\label{sec:finite-gaussian-oracle}

A related oracle inequality holds for an actual finite-scale posterior,
without sending a prior parameter to zero.  The proof below applies to
all learner widths at once when the prior variances decay sufficiently
quickly with width.  It controls clipped posterior draws and the mean
of clipped outputs.  Its likelihood still uses the raw real-valued
network output.

Use the network architecture, energy $Q_N$, and risks from
Section~\ref{sec:gaussian-oracle}.  In this section inputs lie in
$[-1,1]^d$ and labels lie in $[-1,1]$; their joint law is otherwise
arbitrary.  For $\gamma\geq2$---and also for $\gamma\geq1$ in the
one-hidden-layer result below---let the prior $\mu_{N,\gamma}$ have
independent centered Gaussian coordinates with variances
\begin{gather}
 \operatorname{Var}(W^1_{jk})=\frac{N^{1-\gamma}}d,\qquad
 \operatorname{Var}(v_j)=N^{-\gamma-1},\notag\\
 \operatorname{Var}(W^\ell_{jk})=N^{-\gamma}\quad(2\leq\ell\leq H).
 \label{eq:finite-prior}
\end{gather}
Thus this is the Gaussian family of~\eqref{eq:gaussian-prior-variances}
with $\sigma^2=N^{-\gamma}$.  For $\lambda>0$, use a Gaussian likelihood
with variance $n\lambda N^{-\gamma}$ and mean $z_w(X_i)$.  The posterior
is
\begin{equation}
 \Pi_{N,\gamma,S}(dw)
   =\frac{\exp[-\beta_N\widehat R_S(z_w)]}{Z_{N,\gamma,S}}
          \mu_{N,\gamma}(dw),
 \qquad \beta_N=\frac{N^\gamma}{2\lambda}.
 \label{eq:finite-posterior}
\end{equation}
All statements below concern this probability measure at its specified
finite width and scale.

\subsection{An explicit finite-scale oracle inequality}

Let $g$ be a teacher of hidden width $m$, with weights bounded by
$B\geq1$.  Define
\begin{gather}
 E=d(d+1)B^2+2Hm^2B^2,\qquad U=mB,\notag\\
 D_0=d+H+1,\qquad A_g=\sqrt2mB,\notag\\
 K_g=1+A_g\left(\sqrt2 A_g^{H-1}
                       +\sum_{j=0}^{H-2}A_g^j\right).
 \label{eq:finite-teacher-constants}
\end{gather}
The sum is empty if $H=1$.  In particular, $K_g\geq1$ and
$\log K_g=O(H\log(2+mB))$.
For $H\geq2$, put
\begin{gather}
 p=H/2,\qquad
 A_n=\frac{8D_n}{H^p\sqrt n},\qquad
 J_\lambda=\left\lceil\log_2\frac8\lambda\right\rceil,
 \notag\\
 t_\lambda=\log\frac{2(J_\lambda+1)}\delta,
 \qquad
 r_{n,\lambda}=A_n+4\sqrt{\frac{t_\lambda}{2n}}.
 \label{eq:finite-statistical-constants}
\end{gather}
For the scalar teacher concentration term, write
\begin{equation}
 b_n(g,\delta)
   =(U+1)\sqrt{\frac{2R(g)\log(2/\delta)}n}
       +\frac{2(U+1)^2\log(2/\delta)}{3n}.
 \label{eq:finite-teacher-deviation}
\end{equation}

\begin{theorem}[A finite-scale oracle inequality, uniform over widths]
\label{thm:finite-gaussian-oracle}
Let $H\geq2$, $n\geq2$, $0<\delta<1$, and suppose
\begin{equation}
 \frac1n\leq\lambda\leq1,
 \qquad \lambda^p\geq2^{4p-2}A_n.
 \label{eq:finite-calibration}
\end{equation}
For every fixed teacher $g$, with probability at least $1-\delta$,
simultaneously for every $N\geq m$ and every $\gamma\geq2$,
\begin{align}
 \int R(\operatorname{clip}z_w)\,\Pi_{N,\gamma,S}(dw)
 &\leq R(g)+b_n(g,\delta)+r_{n,\lambda}\notag\\
 &\quad+\lambda\left[
 E+1+D_0\log2
       +2D_0\log\frac{4K_g(U+1)}\lambda\right].
 \label{eq:finite-oracle}
\end{align}
The same right-hand side bounds the risk of the posterior mean of
clipped outputs,
\[
 x\longmapsto\int\operatorname{clip}(z_w(x))\,
\Pi_{N,\gamma,S}(dw).
\]

For example, take
\begin{equation}
 \lambda_n=\max\left\{\frac1n,
                 \bigl(2^{4p-2}A_n\bigr)^{1/p}\right\},
 \label{eq:finite-schedule}
\end{equation}
for sufficiently large $n$ that $\lambda_n\leq1$.  At fixed
$d,H,\delta$ and teacher, the excess in~\eqref{eq:finite-oracle} is
$\widetilde O(n^{-1/H})$.  The bound has no learner-width factor and
requires no limit in either $N$ or the prior scale.
\end{theorem}

\subsection{A deviation inequality valid at every energy}

The first step strengthens the deviation bound used for the
strong-prior minimizers: here it must hold over the full parameter
space because the finite-scale posterior has unbounded support.

\begin{lemma}[Global clipped deviation]
\label{lem:finite-global-deviation}
Under~\eqref{eq:finite-calibration}, with probability at least
$1-\delta/2$, simultaneously for every network at every width,
\begin{equation}
 R(\operatorname{clip}z_w)
     \leq\widehat R_S(z_w)+\frac\lambda2 Q_N(w)+r_{n,\lambda}.
 \label{eq:finite-global-deviation}
\end{equation}
\end{lemma}
\begin{proof}
Apply~\eqref{eq:gaussian-fixed-energy-deviation} to the deterministic
classes $\mathcal Z_{2^j}$, $0\leq j\leq J_\lambda$, with individual
failure probability $\delta/[2(J_\lambda+1)]$.
On the intersection event, write $k=Q_N(w)$.  For $k\leq1$, the
unit-energy class gives the claim, since clipping cannot increase
empirical square loss.  For $1<k\leq8/\lambda$, the appropriate shell
has radius at most $2k$, so its capacity term is at most
\[
 A_n2^p k^p
 \leq A_n2^p(8/\lambda)^{p-1}k
 \leq\frac\lambda2 k,
\]
where the last inequality is exactly
$\lambda^p\geq2^{4p-2}A_n$.
Finally, for $k>8/\lambda$, the clipped population loss is at most
four and the empirical loss is nonnegative; therefore
\[
 R(\operatorname{clip}z_w)-\widehat R_S(z_w)
   \leq4<\frac\lambda2 k.
\]
This also covers arbitrarily large energies.
\end{proof}

\subsection{The posterior's variational comparison}

There are
\[
 P_N=N(d+2)+(H-1)N^2
\]
parameters.  Under the Gaussian prior, $N^\gamma Q_N$ has a
chi-square distribution with $P_N$ degrees of freedom.  Thus
\[
 \log\int\exp\left(\frac{N^\gamma Q_N}4\right)d\mu_{N,\gamma}
       =\frac{P_N}2\log2.
\]
The entropy variational inequality implies, for every comparison
measure $Q$,
\begin{equation}
 \operatorname{KL}(Q\Vert\mu_{N,\gamma})
   \geq\frac{N^\gamma}4\int Q_N\,dQ-\frac{P_N}2\log2.
 \label{eq:finite-kl-energy}
\end{equation}
The posterior~\eqref{eq:finite-posterior} minimizes
\[
 Q\longmapsto\int\widehat R_S(z_w)\,dQ
                   +2\lambda N^{-\gamma}
                          \operatorname{KL}(Q\Vert\mu_{N,\gamma}).
\]
Combining its variational optimality with~\eqref{eq:finite-kl-energy}
gives, for every comparison measure $Q$ of finite relative entropy,
\begin{align}
 &\int\left[\widehat R_S(z_w)+\frac\lambda2 Q_N(w)\right]
                          \Pi_{N,\gamma,S}(dw)\notag\\
 &\hspace{8mm}\leq
      \int\widehat R_S(z_w)\,Q(dw)
       +2\lambda N^{-\gamma}\operatorname{KL}(Q\Vert\mu_{N,\gamma})
       +\lambda\frac{P_N}{N^\gamma}\log2.
 \label{eq:finite-variational-comparison}
\end{align}
All quantities are finite for the Gaussian comparison measure used
below.  They are also finite for the posterior: its density relative
to the prior is bounded, and the prior has finite energy moments.
For $\gamma\geq2$ and $N\geq1$,
\begin{equation}
 \frac{P_N}{N^\gamma}\leq d+2+H-1=D_0.
 \label{eq:finite-dimension-ratio}
\end{equation}
For $H\geq2$, this is the point at which the condition
$\gamma\geq2$ enters.  When $H=1$, the quadratic term in $P_N$ is
absent: $P_N/N^\gamma\leq d+2=D_0$ already holds for every
$\gamma\geq1$.

\subsection{A Gaussian comparison measure near the teacher}

Let $w_*$ be the exact teacher clone from
Lemma~\ref{lem:gaussian-cloning}.  Denote the prior covariance by
$\Sigma_N$.  For $0<r\leq1$, define
\[
 Q_r=\mathcal N(w_*,r^2\Sigma_N).
\]
The elementary Gaussian relative-entropy formula gives
\begin{equation}
 \operatorname{KL}(Q_r\Vert\mu_{N,\gamma})
  =\frac{N^\gamma Q_N(w_*)}2
       +\frac{P_N}2\bigl(r^2-1-\log r^2\bigr)
  \leq\frac{N^\gamma E}2+P_N\log(1/r).
 \label{eq:finite-comparison-kl}
\end{equation}

We now bound the comparison measure's prediction error.  The clone's
internal matrices have operator norm at most $A_g=\sqrt2mB$, and
$\sqrt N\|v_*\|_2\leq A_g$.  For the internal matrices, identify each
group of clones with its normalized indicator vector.  The induced
matrix has entries $T_{ji}\sqrt{k_j/k_i}$, whose Frobenius norm is at
most $\sqrt2\,\|T\|_F\leq\sqrt2mB$; the readout bound follows directly
from the cloning construction.
For fixed $x\in[-1,1]^d$, put
\[
 e_\ell(x)
   =\left(\frac1N\mathbb E_{Q_r}
        \|h^\ell(x)-h_*^\ell(x)\|_2^2\right)^{1/2}.
\]
The activation is $1$-Lipschitz.  In the first layer,
\begin{equation}
 e_1(x)
  \leq rN^{(1-\gamma)/2}\frac{\|(1,x)\|_2}{\sqrt d}
  \leq\sqrt2\,r.
 \label{eq:finite-first-perturbation}
\end{equation}
For later hidden layers, use
\[
 W^\ell h^{\ell-1}-W_*^\ell h_*^{\ell-1}
   =W_*^\ell(h^{\ell-1}-h_*^{\ell-1})
                   +(W^\ell-W_*^\ell)h^{\ell-1}.
\]
The noise matrix in the second term is independent of
$h^{\ell-1}$, and $\|h^{\ell-1}\|_2^2\leq N$.  Hence Minkowski's
inequality and its Gaussian second moment give
\begin{equation}
 e_\ell(x)
   \leq A_g e_{\ell-1}(x)+rN^{(1-\gamma)/2}
   \leq A_g e_{\ell-1}(x)+r.
 \label{eq:finite-hidden-perturbation}
\end{equation}
Readout noise is likewise independent of the last hidden layer and
has prediction variance at most $r^2N^{-\gamma}$.  Consequently
\begin{align}
 \left(\mathbb E_{Q_r}|z_w(x)-g(x)|^2\right)^{1/2}
 &\leq A_g e_H(x)+rN^{-\gamma/2}\notag\\
 &\leq r\left[1+A_g\left(\sqrt2 A_g^{H-1}
                       +\sum_{j=0}^{H-2}A_g^j\right)\right]
   =rK_g.
 \label{eq:finite-output-perturbation}
\end{align}
This is a uniform \emph{pointwise second-moment} bound.  It does not
assert that one draw is uniformly close to the teacher over all inputs;
that stronger assertion is unnecessary here.

Choose
\begin{equation}
 \rho=\frac\lambda{4(U+1)},\qquad r=\frac\rho{K_g}.
 \label{eq:finite-comparison-radius}
\end{equation}
These satisfy $0<r\leq1$.  Minkowski's inequality in the product of
$Q_r$ and the empirical input measure gives
\begin{align}
 \int\widehat R_S(z_w)\,Q_r(dw)
 &\leq\widehat R_S(g)+2\rho\sqrt{\widehat R_S(g)}+\rho^2
 \notag\\
 &\leq\widehat R_S(g)+2\rho(U+1)+\rho^2
 \leq\widehat R_S(g)+\lambda.
 \label{eq:finite-comparison-risk}
\end{align}

\begin{proof}[Proof of Theorem~\ref{thm:finite-gaussian-oracle}]
Integrate the global deviation inequality
\eqref{eq:finite-global-deviation} against the posterior, and then
apply~\eqref{eq:finite-variational-comparison} with $Q=Q_r$.
Equations~\eqref{eq:finite-dimension-ratio},
\eqref{eq:finite-comparison-kl}, and
\eqref{eq:finite-comparison-risk} imply
\begin{align*}
 \int R(\operatorname{clip}z_w)\,\Pi_{N,\gamma,S}(dw)
 &\leq\widehat R_S(g)+r_{n,\lambda}\\
 &\quad+\lambda\left[E+1+D_0\log2
               +2D_0\log\frac{4K_g(U+1)}\lambda\right].
\end{align*}
For the fixed teacher, Bernstein's inequality gives
$\widehat R_S(g)\leq R(g)+b_n(g,\delta)$ except on an event of
probability $\delta/2$.  Intersecting this event with the global
deviation event proves~\eqref{eq:finite-oracle}.  Neither event depends
on $N$ or $\gamma$, so the conclusion is simultaneous over all the
claimed widths and scales.  Jensen's inequality proves the claim for
the mean of clipped outputs.  The rate follows from
\eqref{eq:finite-schedule}; $\log(1/\lambda_n)$ adds a logarithm,
without changing its power of $n$.
\end{proof}

\subsection{One hidden layer and the limits of this argument}

At $H=1$, the finite-scale result holds for \emph{every}
$\gamma\geq1$, including the mean-field scale $\gamma=1$, with rate
$\widetilde O(n^{-1/2})$.  Its guarantee is simultaneous over all
$N\geq m$ and all these $\gamma$.  The wider range follows because
$P_N=(d+2)N$, so $P_N/N^\gamma\leq d+2$ for $\gamma\geq1$.
The Gaussian perturbation estimates also hold throughout this range:
$N^{(1-\gamma)/2}\leq1$ and $N^{-\gamma/2}\leq1$ are their only
scale-dependent requirements.  In
\eqref{eq:finite-statistical-constants} take $p=1/2$ and replace the
residual by
\[
 r_{n,\lambda}=A_n+\frac{A_n^2}{\lambda}
                        +4\sqrt{\frac{t_\lambda}{2n}}.
\]
There is no condition $\lambda^p\geq2^{4p-2}A_n$ in this case.
The needed shell estimate is simply
\[
 A_n\sqrt{2k}\leq\frac\lambda2 k+\frac{A_n^2}{\lambda}.
\]
The unit-energy and high-energy cases are unchanged, and all subsequent
steps of the proof hold verbatim.  Choosing
$\lambda_n=\max\{A_n,1/n\}$, for sufficiently large $n$ that
$\lambda_n\leq1$, gives the stated rate for a fixed teacher.

The finite-scale theorem proves a statement about clipped draws and
means of clipped outputs.  It makes no claim about the raw posterior
mean at finite scale.  For two or more hidden layers, this proof does not cover the entire
range $\gamma\geq1$: when $\gamma<2$, the quadratic parameter-count
term makes $P_N/N^\gamma$ grow with the learner width.  Removing that
term requires additional control in function space.  With one hidden
layer, that quadratic term is absent, and the theorem already covers
every $\gamma\geq1$.  The raw-output results
of Section~\ref{sec:raw-gaussian-oracle} instead concern the
strong-prior limit at each fixed width and therefore do not need
$\gamma\geq2$.
